\documentclass[10pt]{article}
\usepackage[margin=0.78in]{geometry}
\usepackage[T1]{fontenc}
\usepackage{lmodern}
\usepackage{amsmath,amssymb}
\usepackage[hidelinks]{hyperref}

\begin{document}
\begin{center}
{\Large\bfseries Literature Status for arXiv:math/0201144:\\
The Little H\"older Space Is an $M$-Ideal\par}
\smallskip
{\small Run fa\_banach\_001, agent\_lane\_05 --- 11 August 2026}
\end{center}
\vspace{-0.7em}

\section*{Status}

This is a \texttt{literature\_already\_answered} packet.  Berninger and
Werner ask whether the pointed little H\"older space
\(H_\alpha^0=\operatorname{lip}_0([0,1],|\cdot|^\alpha)\) is an
\(M\)-ideal in
\(H_\alpha=\operatorname{Lip}_0([0,1],|\cdot|^\alpha)\), for
\(0<\alpha<1\) \cite{BW03}.  Kalton's 2004 Theorems~6.1 and~6.6 and the
remark immediately after Theorem~6.6 answer this affirmatively, in fact for
every compact metric space equipped with a nontrivial gauge \cite{Kalton04}.

\section*{Source Question}

On page~4 of the journal article (the final paragraph of the introduction),
Berninger and Werner state that it remains open whether \(H_\alpha^0\) is an
\(M\)-ideal in \(H_\alpha\).  Section~4 revisits the same question and asks
whether their ``almond function'' might provide a counterexample.  Their
notation is pointed: the functions vanish at the base point \(0\), and the
norm is the H\"older/Lipschitz seminorm, which is then a norm.

\section*{Supporting Answer}

Let \(M\) be compact and let \(\omega\) be a nontrivial gauge.  Kalton's
Theorem~6.1 states
\[
 \operatorname{lip}_\omega(M)^*=\mathcal F_\omega(M),
 \qquad
 \operatorname{lip}_\omega(M)^{**}=\operatorname{Lip}_\omega(M).
\]
Thus \(\operatorname{lip}_\omega(M)\) is the canonical predual of its
Lipschitz-free space.  Theorem~6.6 proves that, for every compact metric
space \(M\) and every \(\varepsilon>0\), \(\operatorname{lip}(M)\) is
\((1+\varepsilon)\)-isometric to a subspace of \(c_0\).  The remark
immediately following the theorem records the relevant consequence:
whenever \(M\) is compact and \(\operatorname{lip}(M)\) is a predual of
\(\mathcal F(M)\), it is an \(M\)-ideal in \(\operatorname{Lip}(M)\).

Apply this with \(M=[0,1]\) and \(\omega(t)=t^\alpha\).  This is a
nontrivial gauge for \(0<\alpha<1\), so Theorem~6.1 supplies the predual
hypothesis and the post-Theorem~6.6 remark gives
\[
 H_\alpha^0
 =\operatorname{lip}_\omega([0,1])
 \quad\text{is an }M\text{-ideal in}\quad
 H_\alpha=\operatorname{Lip}_\omega([0,1]).
\]
Hence the answer to the source question is \emph{yes}.

\section*{Scope and provenance}

The answer is stronger than the interval problem: the same implication holds
for every compact metric \(M\) with a nontrivial gauge.  The identification is
not a new proof; it is a direct application explicitly signposted by Kalton's
remark.  A later thesis and conference abstract by Niglas independently name
Kalton's Theorem~6.6 as the resolution of the Berninger--Werner problem.

{\small
\begin{thebibliography}{9}

\bibitem{BW03}
H.~Berninger and D.~Werner,
\emph{Lipschitz spaces and $M$-ideals},
Extracta Math. \textbf{18} (2003), 33--56;
arXiv:math/0201144.

\bibitem{Kalton04}
N.~J. Kalton,
\emph{Spaces of Lipschitz and H\"older functions and their applications},
Collect. Math. \textbf{55} (2004), 171--217,
\href{https://doi.org/10.1344/CM.V55I2.4055}{doi:10.1344/CM.V55I2.4055}.

\end{thebibliography}
}

\end{document}
